\section{Fazit}

Die zentrale Motivation dieses Projekts hat sich mit der Analyse -- durch Methoden der Digital Humanities -- des filmischen Sehverhaltens von Personen beschäftigt.
Die Methodik sollte primär Auskunft über eine mögliche Clusterbildung zwischen den Nutzern geben und weiterführend einzelne Faktoren hinsichtlich ihrer Auswirkung auf die Gruppenbildung und möglicher Parallelen untersuchen.

Die Netzwerkanalyse hat gezeigt, dass es eine klar definierte Gruppenunterteilung zwischen den Zuschauern gibt.
Durch die Abbildung 1b wurde deutlich gemacht, dass vereinzelt Gruppen existieren, welche in direkter Korrelation zu einem oder mehreren Genres stehen.
Dies ist allerdings nicht die Norm.
Weiterhin hat die genaue Betrachtung der für jede Gruppe dominierenden Filmart eine klare Tendenz zu den Blockbusterfilmen deutlich gemacht, während keine Gruppen stellvertretend für Nischenfilme gefunden werden konnten.
Es wurde folgend in 7.2 gezeigt, wie die Zeit als Maß Auswirkungen auf die Clusterbildung hat.
Während in Abbildung 4a noch klar voneinander visuell trennbare Cluster zu erkennen waren, stellt der moderne Graph (Abbildung 3a) eine klare Verschwimmung der Communitys dar.

Die Methodik der Netzwerkanalyse ist für die technische Auswertung der Daten mit Blick auf die gestellte Forschungsfrage perfekt geeignet gewesen.
Der Informationsverlust bezüglich der Gemeinsamkeitsmetrik zweier Personen erschwert es jedoch, direkte Filmempfehlungen zu geben.
Auch wenn zwei Nutzer ein in dem Graphen stark ähnliches Sehverhalten aufzeigen, ist nicht direkt ersichtlich, welche Faktoren dafür von signifikanter Relevanz waren.

Für zukünftige Arbeiten, die auf einer ähnlichen Forschungsfrage basieren, sollte ein größeres Datenvolumen analysiert werden, um vielleicht existierende Nischengruppen, welche in diesem Projekt nicht gefunden wurden, zu erkennen.
Weiterhin sind die unterschiedlichen Reviewzeiten jedes Nutzers problematisch.
Die wenigsten Nutzer sind über den Gesamtzeitraum von 1995 bis 2009 aktiv auf der untersuchten Plattform gewesen, wodurch Inkonsistenzen in ihrem Sehverhalten entstehen können.
Abschließend könnte neben der puren Bewertung auch die geschriebene Rezension selbst als weiterer Faktor dienen, um das Sehverhalten noch genauer beschreiben zu können.




